\documentclass[11pt]{article}
\usepackage[margin=2.1cm]{geometry}
\usepackage{booktabs}
\usepackage{amsmath}
\usepackage{amssymb}
\usepackage{enumitem}
\usepackage{array}
\usepackage{listings}
\usepackage{xcolor}
\usepackage[T1]{fontenc}
\usepackage{lmodern}

\setlength{\parindent}{0pt}
\setlength{\parskip}{3.6pt}
\newcommand{\code}[1]{\texttt{#1}}
\newcommand{\hr}{\vspace{2pt}\hrule\vspace{6pt}}
\renewcommand{\arraystretch}{1.12}

\definecolor{kw}{RGB}{0,90,160}
\definecolor{cmt}{RGB}{110,110,110}
\lstdefinestyle{vhdl}{
  language=VHDL, basicstyle=\ttfamily\footnotesize,
  keywordstyle=\color{kw}\bfseries, commentstyle=\color{cmt}\itshape,
  showstringspaces=false, columns=fullflexible, breaklines=true,
  frame=leftline, framesep=6pt, xleftmargin=8pt, aboveskip=6pt, belowskip=6pt
}
\lstset{style=vhdl}

\begin{document}

\begin{center}
{\Large\bfseries EE\,362 --- Pre-Lab 7: Solution}\\[2pt]
{\large Next Instruction Address Generation I}\\[4pt]
Fall 2026 \quad$\cdot$\quad 5 points \quad$\cdot$\quad TA marking guide
\end{center}
\hr

{\small Marking: Q1 = 2 points, Q2 = 1 point, Q3 = 2 points. Q1 is pure
arithmetic --- mark it exactly, because every wrong answer here becomes a CPU
that jumps to the wrong place in Lab 12.}

\vspace{4pt}\hr

\section*{Question 1 --- Branch arithmetic (2 points)}

\subsection*{The three steps, every time}

\[
\code{NSA} = PC + 4
\qquad
\text{offset} = \text{sign-extend}(\code{Immed16}) \times 4
\qquad
\text{target} = \code{NSA} + \text{offset}
\]

\subsection*{(a) \code{BLEZ r1, +2} at \code{0x0000000C}, encoding \code{0x18200002}}

\begin{center}
\begin{tabular}{lll}
\toprule
step & value & working \\
\midrule
\code{NSA} & \code{0x00000010} & $\code{0x0C} + 4$ \\
offset & \code{0x00000008} & \code{Immed16} $=$ \code{0x0002}, positive, $2 \times 4 = 8$ \\
\textbf{target} & \textbf{\code{0x00000018}} & $\code{0x10} + \code{0x08}$ \\
\bottomrule
\end{tabular}
\end{center}

\subsection*{(b) Backward branch: $PC = \code{0x30}$, \code{Immed16} $=$ \code{0xFFFC}}

\begin{center}
\begin{tabular}{lll}
\toprule
step & value & working \\
\midrule
\code{NSA} & \code{0x00000034} & $\code{0x30} + 4$ \\
offset & \code{0xFFFFFFF0} & \code{0xFFFC} sign-extends to $-4$; $-4 \times 4 = -16$ \\
\textbf{target} & \textbf{\code{0x00000024}} & $\code{0x34} - \code{0x10}$ \\
\bottomrule
\end{tabular}
\end{center}

\textbf{The teaching point:} nothing special was done for the negative case. The
offset was sign-extended, shifted, and \emph{added}. \textbf{Two's complement
addition handles backward branches with no extra hardware} --- there is no
subtractor anywhere in the branch path.

\subsection*{(c) The two mistakes}

\begin{center}
\begin{tabular}{lll}
\toprule
mistake & computed target & what it looks like \\
\midrule
forgot the $\times 4$ shift & \code{0x10} $+$ \code{2} $=$ \textbf{\code{0x12}} & an \textbf{unaligned} address \\
used $PC$ instead of $PC+4$ & \code{0x0C} $+$ \code{8} $=$ \textbf{\code{0x14}} & a \textbf{valid} instruction address, off by one instruction \\
\bottomrule
\end{tabular}
\end{center}

\begin{quote}
\textbf{Model answer for ``which is harder to notice'':} the $PC$-versus-$PC+4$
mistake. It produces a properly aligned address that points at a real
instruction, so the machine keeps running and simply executes the wrong one ---
the program computes a wrong answer with no obvious symptom. The missing shift
gives \code{0x12}, which is not a multiple of four; the instruction memory index
is then wrong in an obvious way and the failure is immediate and loud.
\end{quote}

Full credit needs \emph{both} numbers and the argument that a plausible-looking
address is more dangerous than an impossible one.

\vspace{4pt}\hr

\section*{Question 2 --- Interface and reuse (1 point)}

\subsection*{(a) The entity}

\begin{lstlisting}
entity NIAG_Reg is
    generic (
        RESET_VEC : natural := 0          -- where the machine starts
    );
    port (
        NextPC : in  std_logic_vector(31 downto 0);
        Strb   : in  std_logic;           -- load enable
        Rst    : in  std_logic;           -- asynchronous
        clk    : in  std_logic;
        PC     : out std_logic_vector(31 downto 0)
    );
end entity NIAG_Reg;
\end{lstlisting}

The generic must be there, and it must be usable as a \emph{value}, not a width.
That is the only structural difference from \code{RegN}.

\subsection*{(b) The equation, and what changed}

\[
PC^{+} = \overline{\code{Rst}} \cdot \left( \code{Strb} \cdot \code{NextPC}
   + \overline{\code{Strb}} \cdot PC \right)
\]

\begin{quote}
\textbf{Model answer:} It is exactly Lab 2's register equation with
\code{Ld} renamed \code{Strb}. \textbf{The only real change is the reset value:}
\code{RegN} resets to zero, the program counter resets to the address where the
machine starts.
\end{quote}

Full credit needs ``same equation, different reset value''. A student who writes
$PC^{+} = PC + 4$ has confused the \textbf{register} with the \textbf{adder} that
feeds it --- worth correcting immediately, because keeping those two separate is
the whole structure of this lab.

\vspace{4pt}\hr

\section*{Question 3 --- Test plan (2 points)}

\subsection*{Model vector list}

\begin{center}
\small
\begin{tabular}{c>{\raggedright\arraybackslash}p{5.2cm}>{\raggedright\arraybackslash}p{7.4cm}}
\toprule
\# & inputs & what it proves \\
\midrule
1 & \code{Rst} = 1 & the PC loads the reset vector, not zero (unless they coincide) \\
2 & \code{Strb} = 1, one edge & the PC takes \code{NextPC} \\
3 & \code{Strb} = 0, one edge & the PC \textbf{holds}: an edge alone does not advance it \\
4 & \code{PC} = \code{0x0C}, read \code{NSA} & $\code{NSA} = PC+4$, and it is \textbf{combinational} \\
5 & forward branch, \code{Immed16} = \code{0x0002} & offset shifted by 2 and added to \code{NSA} \\
6 & \textbf{backward} branch, \code{Immed16} = \code{0xFFFC} & \textbf{sign extension works}: target is \emph{below} \code{NSA} \\
7 & \code{Immed16} = \code{0x0000} & the degenerate branch: target $=$ \code{NSA} \\
\bottomrule
\end{tabular}
\end{center}

Six vectors earn full credit. \textbf{Vector 6 is the one that must appear.}

\subsection*{(a) If every offset is positive}

\begin{quote}
\textbf{Model answer:} You could not detect a \textbf{zero-extended} immediate
where sign extension was required. For positive offsets the two are identical, so
every vector passes; the design then jumps forward by an enormous amount on the
first backward branch.
\end{quote}

\textbf{This is the same lesson as Lab 4:} a test value must be able to
distinguish the cases. And it matters here more than anywhere, because
\textbf{every loop in the Lab 12 program ends with a backward branch} --- a CPU
with this bug runs straight past the end of its loop.

\subsection*{(b) Proving \code{NSA} is combinational}

\begin{quote}
\textbf{Model answer:} Change \code{PC} \emph{between} clock edges and check that
\code{NSA} changes immediately, with no edge in between --- for example, drive
\code{PC} to \code{0x0C}, wait, observe \code{0x10}; then drive \code{0x10}, wait,
observe \code{0x14}, all within one clock phase.
\end{quote}

\begin{quote}
\textbf{What would go wrong if \code{NSA} were clocked:} the next address would
arrive one cycle late. The CPU needs \code{NSA} \emph{during} the same cycle in
which it fetches, because the adder's result must be ready at the edge that loads
the PC. A registered \code{NSA} would make the machine compute the next address
from the \emph{previous} PC --- it would execute every instruction twice, or skip
every other one, depending on the timing.
\end{quote}

\textbf{The structural rule behind the whole lab:} the program counter is the only
storage element here. \code{NSA} and \code{Rel} are pure arithmetic and must stay
combinational.

\vspace{6pt}
\hr
{\small \textbf{At the bench:} check Q1(b) first. A student whose backward target
is above \code{NSA}, or who wrote \code{0x44}, has either dropped the sign
extension or added instead of letting two's complement do the work.}

\end{document}
